% Figure: interference (shrink) fit of a hub on a solid shaft. Cross-
% section sketch on the left shows the solid shaft, the hub bore, and the
% radial interference delta (exaggerated) that the assembly squeezes out,
% generating a uniform contact pressure p at the interface. The pressure
% drives a tensile hoop stress in the hub (largest at the bore) and a
% compressive radial squeeze on the shaft. Right side names the stress
% directions at the interface with short arrows.
\begin{figure}[htbp]
\centering
\begin{tikzpicture}[scale=1.0]
% left: half cross-section
% shaft (solid), radius a = 1.4 cm drawn
\fill[enggray!25] (0,-1.4) rectangle (2.2,1.4);
\draw[engblack, very thick] (0,-1.4) -- (2.2,-1.4);
\draw[engblack, very thick] (0,1.4) -- (2.2,1.4);
\node[engblack, font=\scriptsize] at (1.1,0) {shaft};
% hub (outer annulus), from a to b = 2.6 cm
\fill[engblue!12] (0,1.4) rectangle (2.2,2.6);
\fill[engblue!12] (0,-2.6) rectangle (2.2,-1.4);
\draw[engblue, very thick] (0,2.6) -- (2.2,2.6);
\draw[engblue, very thick] (0,-2.6) -- (2.2,-2.6);
\node[engblue, font=\scriptsize] at (1.1,2.0) {hub};
\node[engblue, font=\scriptsize] at (1.1,-2.0) {hub};
% centre line
\draw[enggray, dashed] (-0.1,0) -- (2.4,0);
% interface line (contact surface at radius a)
\draw[engred, very thick] (0,1.4) -- (2.2,1.4);
\draw[engred, very thick] (0,-1.4) -- (2.2,-1.4);
% contact-pressure arrows pushing apart across the interface
\draw[engred, -{Latex}] (0.5,1.25) -- (0.5,1.55);
\draw[engred, -{Latex}] (0.5,1.55) -- (0.5,1.25);
\draw[engred, -{Latex}] (1.5,1.25) -- (1.5,1.55);
\node[engred, font=\scriptsize, anchor=west] at (2.25,1.4) {contact pressure $p$};
% radial dimension markers
\draw[enggray, {Latex}-{Latex}] (-0.05,0) -- (-0.05,1.4);
\node[enggray, font=\scriptsize, anchor=east] at (-0.1,0.7) {$a$};
\draw[enggray, {Latex}-{Latex}] (-0.05,1.4) -- (-0.05,2.6);
\node[enggray, font=\scriptsize, anchor=east] at (-0.1,2.0) {to $b$};
% right: stress directions at the bore of the hub
\begin{scope}[shift={(6.0,0)}]
\node[engblack, font=\small, anchor=south] at (0,2.7) {at the hub bore};
% small element
\draw[engblack, very thick] (-0.7,-0.7) rectangle (0.7,0.7);
% hoop (tensile) arrows, horizontal outward
\draw[engblue, very thick, -{Latex}] (0.7,0) -- (1.5,0);
\draw[engblue, very thick, -{Latex}] (-0.7,0) -- (-1.5,0);
\node[engblue, font=\scriptsize, anchor=west] at (1.5,0.2) {$\sigma_\theta>0$ (tension)};
% radial (compressive) arrows, vertical inward
\draw[engred, very thick, {Latex}-] (0,0.7) -- (0,1.45);
\draw[engred, very thick, {Latex}-] (0,-0.7) -- (0,-1.45);
\node[engred, font=\scriptsize, anchor=south] at (0,1.5) {$\sigma_r=-p$};
\end{scope}
\end{tikzpicture}
\caption[Interference (shrink) fit of a hub on a shaft]{An interference
fit. The hub bore is machined slightly smaller than the shaft, by a
radial interference $\delta$; forcing or heat-shrinking the hub onto the
shaft squeezes that interference out and leaves a uniform contact
pressure $p$ at the interface. The pressure pushes the hub outward,
raising a tensile hoop stress that is largest at the bore, while it
presses radially inward with $\sigma_r = -p$ on both members at the
interface. The contact pressure is what holds the joint: it sets the
friction torque the fit can transmit and the peak hoop stress the hub
must survive.}
\label{fig:vol03ch03:shrink-fit-pressure}
\end{figure}
